\section{Fog, exposure, tone mapping and gamma}
\label{sec:post}
\here{9}
Lighting produces linear high-dynamic-range radiance. Before it reaches the screen the image passes through the
post-processing chain of \code{ShadersPost.cpp} and \code{ShadersPostFx.cpp}, every step a shader of the project:
\begin{itemize}
\item ambient occlusion (SSAO / GTAO) $\to$ lighting $\to$ combine: volumetric fog, reflections, aerial perspective $\to$ TAA
$\to$ frame statistics (auto-exposure) $\to$ bloom $\to$ composite: exposure, grade, tone curve, standard red, green, blue colour space (sRGB) encoding.
\end{itemize}

\textbf{Fog.} Each surface is mixed towards the fog colour by a factor that grows with distance $d$ from the eye
(\code{ShaderSources.cpp: fogAmount}); the volumetric pass integrates light scattered by the air along the view ray
(\code{ShadersPost.cpp: kFogFrag}):
\begin{gather}
\vect c'=\operatorname{lerp}(\vect c,\ \vect c_{\text{fog}},\ f),\qquad f=1-e^{-(\rho d)^2},\quad \rho=0.018\ \text{(night)}\ \ldots\ 0.0042\,\text{m}^{-1}\ \text{(day)}\label{eq:fog}\\
\sigma(\vect x)=\sigma_0\,e^{-0.22\,y}\big(0.45+1.1\,\text{noise}\big),\qquad T_{i+1}=T_i\,e^{-\sigma_i\Delta s}\nonumber\\
\vect L_{\text{in}}\mathrel{+}=T_i\,\vect L_{\text{light}}\,p_{\text{HG}}(\cos\theta)\,\big(1-e^{-\sigma_i\Delta s}\big),\qquad
\vect c'=\vect c\,T+\vect L_{\text{in}}\nonumber
\end{gather}
\begin{itemize}
\item The squared exponent ($\exp^2$ fog) keeps the near field clear and closes in faster than plain $\exp$ fog
(Figure~\ref{fig:fogcurve}); $\rho$ is interpolated between the night and day values by the solar model.
\item The volumetric march takes 24--64 steps (preset) with a per-pixel jitter; $p_{\text{HG}}(c)=\frac{1-g^2}{4\pi(1+g^2-2gc)^{3/2}}$
is the Henyey--Greenstein phase function ($g=0.55$ moon, $0.35$ searchlight); the searchlight shaft is marched only
inside its cone (an analytic ray--cone interval).
\end{itemize}

\textbf{Exposure and tone mapping} (\code{ShadersPostFx.cpp: kCompositeFrag}, CPU mirror \code{math/ToneMap.cpp}):
\begin{gather}
k=k_0\,2^{\operatorname{clamp}(\kappa-\overline{\log_2 Y},\,-1.5,\,1.5)},\qquad Y=0.2126R+0.7152G+0.0722B\label{eq:tonemap}\\
\text{ACES fit: } \vect x\mapsto M_{\text{out}}\,\frac{\vect v(\vect v+0.0245786)-0.000090537}{\vect v(0.983729\,\vect v+0.432951)+0.238081},\qquad \vect v=M_{\text{in}}\vect x\nonumber\\
\text{AgX: } \vect v=M_{\text{inset}}\vect x,\ \ \ell=\frac{\operatorname{clamp}(\log_2\vect v,-12.47,4.03)+12.47}{16.5},\ \ 
\vect y=\Big(M_{\text{outset}}\sum_{j=0}^{6}a_j\ell^{\,j}\Big)^{2.2}\nonumber
\end{gather}
\begin{itemize}
\item $\overline{\log_2Y}$ is the centre-weighted mean of the frame's log luminance, reduced on the GPU to one texel and
adapted over time; $\kappa$ (the key) follows the time of day.
\item By day the AgX curve (log encoding, a sixth-order sigmoid fit, $a_j$ in \code{ToneMap.h}) desaturates very bright
light towards white; at night the Academy Color Encoding System (ACES) fit; the two cross-fade through twilight. Both compress any radiance into $[0,1]$
(Figure~\ref{fig:tonecurve}a).
\end{itemize}

\textbf{Gamma: linear to sRGB and back.}
\begin{equation}
\text{encode: } c_s=\begin{cases}12.92\,c & c\le0.0031308\\ 1.055\,c^{1/2.4}-0.055 & \text{else}\end{cases},\qquad
\text{decode of a base colour: } \rho=c_s^{\,2.2}\ \ (\code{toLinear})\label{eq:gamma}
\end{equation}
\begin{itemize}
\item The composite pass encodes once, at the very end (\code{linearToSrgb}); the default framebuffer is not an sRGB
framebuffer, so the GPU does no conversion of its own.
\item Material colours are authored in sRGB and decoded with the power 2.2 (close to the exact sRGB curve, Eq.~\eqref{eq:srgb}),
so lighting and blending happen in linear space (Figure~\ref{fig:tonecurve}b).
\end{itemize}

\textbf{Bloom.} Pixels above a threshold $\tau$ glow: a soft knee
$s=\operatorname{clamp}(m-\tau+\tfrac\tau2,0,\tau)^2/(2\tau)$ with $m=\max(R,G,B)$ keeps $\max(s,m-\tau)/m$ of each pixel,
which is then blurred down a chain of six half-size levels (13-tap filter) and back up ($3\times3$ tent, weights
$\frac1{16}[1\,2\,1;2\,4\,2;1\,2\,1]$), and added before exposure.
\begin{figure}[htbp]
\centering\plot{p_fog}{0.5\linewidth}
\caption{The fog factor of Eq.~\eqref{eq:fog} for the night and day densities, against plain exponential fog.}
\label{fig:fogcurve}
\end{figure}
\begin{figure}[htbp]
\centering\plot{p_tonemap}{\linewidth}
\caption{(a) The tone curves of Eq.~\eqref{eq:tonemap} on the grey axis after each one's exposure, with a hard clip for
comparison; (b) the sRGB encoding of Eq.~\eqref{eq:gamma} against a pure power $1/2.2$, and the decoding of base colours.}
\label{fig:tonecurve}
\end{figure}
\begin{figure}[htbp]
\centering
\begin{subfigure}{0.49\linewidth}\raw{lab/explorer_stage_12_fog}{\linewidth}\caption{Fog stage: in-scattered light and
transmittance along each ray.}\end{subfigure}\hfill
\begin{subfigure}{0.49\linewidth}\raw{lab/explorer_stage_14_tonemap}{\linewidth}\caption{Tone-mapping stage: linear clip
(left of the wipe) against the filmic curve, with the luminance histogram.}\end{subfigure}
\caption{The two stages in the Frame Explorer.}
\label{fig:poststages}
\end{figure}
\begin{table}[htbp]
\caption{Post-processing equations and their code.}\label{tab:postcode}
\centering\small
\begin{tabularx}{\linewidth}{@{}lXl@{}}
\toprule
Step & Equation & Code\\\midrule
Distance fog & $f=1-e^{-(\rho d)^2}$ & \code{fogAmount}: \code{1.0 - exp(-pow(uFogDensity * dist, 2.0))}\\
Volumetric fog & Eq.~\eqref{eq:fog}, $T\gets T\,e^{-\sigma\Delta s}$ & \code{kFogFrag}\\
Auto-exposure & Eq.~\eqref{eq:tonemap}, key $\pm1.5$ stops & \code{kReduceFrag}, \code{kCompositeFrag}\\
Tone curve & ACES fit / AgX & \code{acesFitted}, \code{agxDisplay}\\
sRGB encoding & Eq.~\eqref{eq:gamma} & \code{linearToSrgb}\\
Bloom & soft knee, 6 levels & \code{kBloomDownFrag}, \code{kBloomUpFrag}\\
\bottomrule
\end{tabularx}
\end{table}

\textbf{Limitation.} The image is stored in 8 bits per channel after the tone curve; dark gradients can band, and only
the film grain added in the last pass breaks the steps up.
